\section{Condiciones de frontera y coeficientes de Fresnel en dieléctricos}
\label{sec:appendix_maxwell_frontera_fresnel}

En este apéndice se deducen las condiciones de frontera de Maxwell en una interfaz dieléctrica y las expresiones de reflexión/transmisión para ondas planas en polarizaciones $s$ y $p$.

\subsection{Condiciones de frontera en una interfaz dieléctrica}
\label{subsec:apx_condiciones_frontera}

Partimos de las ecuaciones de Maxwell en medios materiales,
\begin{equation}
\begin{aligned}
\nabla \cdot \mathbf{D} &= \rho_f, & \nabla \cdot \mathbf{B} &= 0, \\
\nabla \times \mathbf{E} &= -\frac{\partial \mathbf{B}}{\partial t},
& \nabla \times \mathbf{H} &= \mathbf{J}_f + \frac{\partial \mathbf{D}}{\partial t}.
\end{aligned}
\label{eq:apx_maxwell_materiales}
\end{equation}

Sea la interfaz
\begin{equation}
\Sigma = \{(x,y,z)\mid z=0\},
\end{equation}
que separa dos medios ($z<0$ y $z>0$). Para un campo vectorial genérico,
\begin{equation}
\tilde{\mathbf F}(\mathbf r,t)
=
\mathbf F(\mathbf r,t)\,\Theta(-z)
+
\mathbf F'(\mathbf r,t)\,\Theta(z),
\label{eq:apx_campo_partido}
\end{equation}
con
\begin{equation}
\Theta(z)=
\begin{cases}
1, & z>0,\\
0, & z\le 0.
\end{cases}
\end{equation}

La divergencia y el rotacional toman la forma distribucional:
\begin{align}
\nabla\cdot\tilde{\mathbf F}
&=
\left[\nabla\cdot\mathbf F\right]_H
+
\delta(z)\,\hat{\mathbf n}\cdot\left(\mathbf F'-\mathbf F\right),
\label{eq:apx_div_partida}\\
\nabla\times\tilde{\mathbf F}
&=
\left[\nabla\times\mathbf F\right]_H
+
\delta(z)\,\hat{\mathbf n}\times\left(\mathbf F'-\mathbf F\right).
\label{eq:apx_rot_partida}
\end{align}

Descomponemos fuentes volumétricas y superficiales como
\begin{equation}
\rho_f = \tilde\rho_f + \sigma_f(x,y,t)\,\delta(z),
\qquad
\mathbf J_f = \tilde{\mathbf J}_f + \mathbf K_f(x,y,t)\,\delta(z).
\label{eq:apx_fuentes_partidas}
\end{equation}

Sustituyendo \eqref{eq:apx_div_partida}, \eqref{eq:apx_rot_partida} y \eqref{eq:apx_fuentes_partidas} en \eqref{eq:apx_maxwell_materiales}, y agrupando términos proporcionales a $\delta(z)$, se obtienen las condiciones de frontera:
\begin{equation}
\begin{aligned}
\hat{\mathbf n}\cdot(\mathbf D_2-\mathbf D_1) &= \sigma_f,\\
\hat{\mathbf n}\cdot(\mathbf B_2-\mathbf B_1) &= 0,\\
\hat{\mathbf n}\times(\mathbf E_2-\mathbf E_1) &= 0,\\
\hat{\mathbf n}\times(\mathbf H_2-\mathbf H_1) &= \mathbf K_f.
\end{aligned}
\label{eq:apx_condiciones_frontera_finales}
\end{equation}

\subsection{Ondas planas y relaciones de dispersión}
\label{subsec:apx_ondas_planas}

En ausencia de fuentes libres ($\rho=0$, $\mathbf J=0$):
\begin{equation}
\begin{aligned}
\nabla \cdot \mathbf E &= 0, & \nabla \cdot \mathbf B &= 0,\\
\nabla \times \mathbf E &= -\frac{\partial \mathbf B}{\partial t},
& \nabla \times \mathbf B &= \frac{1}{c^2}\frac{\partial \mathbf E}{\partial t},
\end{aligned}
\label{eq:apx_maxwell_vacio_local}
\end{equation}
con
\begin{equation}
c = \frac{1}{\sqrt{\mu\epsilon}}.
\label{eq:apx_c_medio}
\end{equation}

Para una onda plana armónica,
\begin{equation}
\mathbf E = \mathbf E_0 e^{i(\mathbf k\cdot\mathbf r-\omega t)},
\qquad
\mathbf B = \mathbf B_0 e^{i(\mathbf k\cdot\mathbf r-\omega t)},
\label{eq:apx_ansatz_onda_plana}
\end{equation}
se obtiene
\begin{equation}
\begin{aligned}
i\,\mathbf k\cdot\mathbf E &= 0,
& i\,\mathbf k\cdot\mathbf B &= 0,\\
i\,\mathbf k\times\mathbf E &= i\omega\mathbf B,
& \mathbf k\times\mathbf B &= -\frac{\omega}{c^2}\mathbf E.
\end{aligned}
\label{eq:apx_maxwell_k_omega}
\end{equation}

Aplicando $\mathbf k\times(\mathbf k\times\mathbf E)=\mathbf k(\mathbf k\cdot\mathbf E)-k^2\mathbf E$ y usando \eqref{eq:apx_maxwell_k_omega}, resulta
\begin{equation}
k^2 = \frac{\omega^2}{c^2},
\qquad
\omega = ck,
\qquad
v = \frac{\omega}{k}.
\label{eq:apx_dispersion}
\end{equation}

\subsection{Interfaz plana y coeficientes de Fresnel}
\label{subsec:apx_fresnel}

Consideramos tres ondas planas en la interfaz $\Sigma$: incidente, reflejada y transmitida,
\begin{equation}
\mathbf E_j = \mathbf E_j^0 e^{i(\mathbf k_j\cdot\mathbf r-\omega t)},
\qquad j\in\{i,R,T\}.
\label{eq:apx_campos_iRT}
\end{equation}

De \eqref{eq:apx_condiciones_frontera_finales} se obtiene el sistema base:
\begin{equation}
\begin{aligned}
\bigl(\epsilon(\mathbf E_i^0+\mathbf E_R^0)-\epsilon'\mathbf E_T^0\bigr)\cdot\hat{\mathbf n} &= 0,\\
\bigl[(\mathbf k_i\times\mathbf E_i^0+\mathbf k_R\times\mathbf E_R^0)-(\mathbf k_T\times\mathbf E_T^0)\bigr]\cdot\hat{\mathbf n} &= 0,\\
(\mathbf E_i+\mathbf E_R-\mathbf E_T)\times\hat{\mathbf n} &= 0,\\
\left[\frac{1}{\mu}(\mathbf k_i\times\mathbf E_i^0+\mathbf k_R\times\mathbf E_R^0)-\frac{1}{\mu'}(\mathbf k_T\times\mathbf E_T^0)\right]\times\hat{\mathbf n} &= 0.
\end{aligned}
\label{eq:apx_sistema_frontera_iRT}
\end{equation}

Descomponemos
\begin{equation}
\mathbf E_0 = E_s^0\,\hat{\mathbf s}+E_p^0\,\hat{\mathbf p},
\label{eq:apx_descomposicion_sp}
\end{equation}
y para el modo $s$ se usa además
\begin{equation}
E_i^0+E_R^0=E_T^0,
\label{eq:apx_continuidad_tangencial_s}
\end{equation}
junto con
\begin{equation}
(\mathbf k\times\mathbf E)\times\hat{\mathbf n}
=
(\mathbf k\cdot\hat{\mathbf n})\mathbf E
-
(\hat{\mathbf n}\cdot\mathbf E)\mathbf k.
\label{eq:apx_identidad_triple}
\end{equation}

Definiendo
\begin{equation}
R=\frac{E_R^0}{E_i^0},
\qquad
T=\frac{E_T^0}{E_i^0},
\qquad
Z=\sqrt{\frac{\epsilon}{\mu}},
\qquad
Z'=\sqrt{\frac{\epsilon'}{\mu'}},
\label{eq:apx_def_R_T_Z}
\end{equation}
los coeficientes de Fresnel para polarización $s$ quedan:
\begin{align}
R_s
&=
\frac{Z\cos\theta_i - Z'\cos\theta_T}{Z\cos\theta_i + Z'\cos\theta_T},
\label{eq:apx_fresnel_rs}\\
T_s
&=
\frac{2Z\cos\theta_i}{Z\cos\theta_i + Z'\cos\theta_T}.
\label{eq:apx_fresnel_ts}
\end{align}

Para polarización $p$:
\begin{align}
R_p
&=
\frac{Z\cos\theta_T - Z'\cos\theta_i}{Z\cos\theta_T + Z'\cos\theta_i},
\label{eq:apx_fresnel_rp}\\
T_p
&=
\frac{2Z'\cos\theta_i}{Z\cos\theta_T + Z'\cos\theta_i}.
\label{eq:apx_fresnel_tp}
\end{align}

Las expresiones \eqref{eq:apx_fresnel_rp} y \eqref{eq:apx_fresnel_tp} condensan los resultados finales para el modo $p$ en incidencia oblicua entre dieléctricos.
